% To:
\section{Appendix: Bell Inequalities and the Tsirelson Bound}
\label{app:bell}

This note collects the mathematical background behind the claims made in Sections~4 and~5 regarding Bell's theorem and its consequences. It is intended as a self-contained reference for readers unfamiliar with the formalism.

%% -----------------------------------------------------------
\subsection*{1. The EPR argument and the problem of correlations}

Einstein, Podolsky, and Rosen argued in 1935 \cite{einstein_1935_epr} that if two spatially separated particles produce perfectly anti-correlated measurement outcomes without any possible causal connection between the measurements, then the outcome of each measurement must have been determined in advance by some property carried locally by each particle---a \textit{hidden variable}. Quantum mechanics, they argued, is incomplete because it does not specify this variable.

\medskip

Bell \cite{bell_1964_epr} converted this from a philosophical dispute into a precise mathematical question: what statistical correlations are compatible with \textit{any} local hidden-variable theory, regardless of its internal mechanism?

%% -----------------------------------------------------------
\subsection*{2. Bell's original inequality and Pearson correlations}

Suppose two particles, $A$ and $B$, are produced together and travel to separate detectors. Each detector can be oriented in one of several directions, and each measurement yields an outcome $\pm1$. Let $\lambda$ be a hidden variable (or set of variables) distributed according to some probability density $\rho(\lambda)$, unknown but fixed at the moment the particles are produced. The local hidden-variable assumption is that the outcome at $A$ depends only on the local setting $a$ and on $\lambda$---not on what $B$ does---and symmetrically for $B$:
\begin{equation}
  A = A(a,\lambda) \in \{-1,+1\}, \qquad
  B = B(b,\lambda) \in \{-1,+1\}.
  \label{eq:lhv}
\end{equation}

The correlation between the outcomes at settings $a$ and $b$ is then the \textit{Pearson correlation coefficient} of the two $\pm1$ random variables $A$ and $B$:
\begin{equation}
  E(a,b) = \int \rho(\lambda)\, A(a,\lambda)\, B(b,\lambda)\, \mathrm{d}\lambda.
  \label{eq:pearson}
\end{equation}

This is identical to the standard Pearson product-moment correlation $r_{AB} = \langle AB \rangle / (\sigma_A \sigma_B)$, which here reduces to $\langle AB \rangle$ because $\sigma_A = \sigma_B = 1$ for $\pm1$ variables. $E(a,b) = +1$ means perfect agreement, $E(a,b) = -1$ means perfect anti-correlation, and $E(a,b) = 0$ means statistical independence.

\medskip

Bell showed that for any three measurement directions $a$, $b$, $c$, the correlations defined by~(\ref{eq:pearson}) must satisfy:
\begin{equation}
  |E(a,b) - E(a,c)| \leq 1 + E(b,c).
  \label{eq:bell-supp}
\end{equation}

The proof is pure arithmetic: since $A(b,\lambda)B(b,\lambda) = \pm1$, one has $|A(a,\lambda)B(b,\lambda) - A(a,\lambda)B(c,\lambda)| \leq 1 - A(b,\lambda)B(b,\lambda)A(b,\lambda)B(c,\lambda)$ pointwise in $\lambda$, and integrating over $\rho(\lambda)$ yields (\ref{eq:bell-supp}). No assumption about quantum mechanics enters---the inequality holds for \textit{any} distribution $\rho(\lambda)$ and any $\pm1$-valued functions $A(a,\lambda)$, $B(b,\lambda)$.

%% -----------------------------------------------------------
\subsection*{3. The CHSH inequality}

Clauser, Horne, Shimony, and Holt (1969) \cite{clauser_1969_chsh} reformulated Bell's result in a form better suited to real experiments, which suffer from detector inefficiencies and finite statistics. The four authors' initials give the inequality its name: \textbf{CHSH}.

\medskip

The setup uses two settings per detector---$a$ and $a'$ for particle $A$, and $b$ and $b'$ for particle $B$---yielding four correlation functions. The CHSH quantity $S$ is defined as:

\begin{equation}
  S = E(a,b) - E(a,b') + E(a',b) + E(a',b').
  \label{eq:S}
\end{equation}

The letter $S$ stands for the \textit{sum} (or \textit{score}) of these four Pearson correlations, combined with a specific sign pattern chosen so that any local hidden-variable model is forced into an algebraic contradiction. The CHSH inequality states:
\begin{equation}
  |S| \leq 2.
  \label{eq:chsh-supp}
\end{equation}

The proof is again arithmetic: since each $A(\cdot,\lambda), A'(\cdot,\lambda), B(\cdot,\lambda), B'(\cdot,\lambda) \in \{-1,+1\}$, the combination $AB - AB' + A'B + A'B' = A(B-B') + A'(B+B')$ takes values in which one of $(B - B')$ and $(B + B')$ is always zero and the other is $\pm2$, so the whole expression is $\pm2$ pointwise in $\lambda$, and therefore $|S| \leq 2$ after integrating over $\rho(\lambda)$.

%% -----------------------------------------------------------
\subsection*{4. Quantum mechanics violates the CHSH bound}

For the spin-singlet state $|\Psi^-\rangle = \frac{1}{\sqrt{2}} (|\uparrow\rangle_A|\downarrow\rangle_B - |\downarrow\rangle_A|\uparrow\rangle_B)$, quantum mechanics predicts:
\begin{equation}
  E_{\mathrm{QM}}(a,b) = -\cos(\theta_{ab}),
  \label{eq:eqm}
\end{equation}

where $\theta_{ab}$ is the angle between the measurement directions $a$ and $b$ in the Bloch sphere. Note that (\ref{eq:eqm}) is still a Pearson correlation between $\pm1$ outcomes---the quantum calculation gives the same object as (\ref{eq:pearson}), it just predicts a \textit{different} value for it than any local hidden-variable model can.

\medskip

Choosing the optimal angles $a = 0\circ$, $a' = 90\circ$, $b = 45\circ$, $b' = 135\circ$ (so that each pair of settings subtends $45\circ$):
\begin{equation}
  E(a,b) = E(a',b) = E(a',b') = -\cos 45\circ = -\tfrac{1}{\sqrt{2}}, \qquad  E(a,b') = -\cos 135\circ = +\tfrac{1}{\sqrt{2}},
\end{equation}

giving:
\begin{equation}
  S_{\mathrm{QM}}
  = -\tfrac{1}{\sqrt{2}} - \tfrac{1}{\sqrt{2}} - \tfrac{1}{\sqrt{2}} - \tfrac{1}{\sqrt{2}}
  = -2\sqrt{2} \approx -2.828.
\end{equation}

The magnitude $|S_{\mathrm{QM}}| = 2\sqrt{2} > 2$ violates the CHSH bound~(\ref{eq:chsh-supp}). This violation was confirmed experimentally by Aspect, Grangier, and Roger in 1982 \cite{aspect_1982_bell}, and without loopholes by Hensen et al.\ in 2015 \cite{hensen_2015_loophole}. 

%% -----------------------------------------------------------
\subsection*{5. The Tsirelson bound: why not stronger?}

The CHSH bound for classical correlations is $|S| \leq 2$. Quantum mechanics reaches $|S| = 2\sqrt{2} \approx 2.83$. But why does quantum mechanics stop there? Logically, $|S|$ could be as large as $4$ (if all four Pearson correlations had magnitude 1 simultaneously). Is there a deeper reason why the quantum world is \textit{more} non-local than any classical model but \textit{less} non-local than it could logically be?

\medskip

Boris Tsirelson (also transliterated Cirel'son) answered this in 1980 \cite{tsirelson_1980_bound}: \textit{$2\sqrt{2}$ is the maximum value of $|S|$ achievable by any quantum state and any quantum measurements, in any Hilbert space of any dimension.} This is the \textbf{Tsirelson bound}.

\medskip

To see intuitively why it is $2\sqrt{2}$ and not $4$: the four Pearson correlations $E(a,b)$, $E(a,b')$, $E(a',b)$, $E(a',b')$ cannot all equal $\pm1$ simultaneously in quantum mechanics, because the measurements do not commute. Choosing $a$ and $b$ to maximise $E(a,b)$ forces the other correlations into a configuration that cannot simultaneously be maximised---the non-commutativity of quantum observables introduces a geometric constraint that limits the total score $S$.

\medskip

The algebraic proof uses the operator version of $S$:
\begin{equation}
  \hat{S} = \hat{A}\otimes(\hat{B}+\hat{B}') + \hat{A}'\otimes(\hat{B}-\hat{B}'),
\end{equation}

where $\hat{A}, \hat{A}', \hat{B}, \hat{B}'$ are observables with eigenvalues $\pm1$ (Hermitian operators with $\hat{A}^2 = \mathbb{1}$, etc.). Computing $\hat{S}^2$ and using $\hat{A}^2 = \hat{A}'^2 = \hat{B}^2 = \hat{B}'^2 = \mathbb{1}$:
\begin{equation}
  \hat{S}^2 = 4\,\mathbb{1} + [\hat{A},\hat{A}']\otimes[\hat{B},\hat{B}'].
  \label{eq:S2}
\end{equation}
The operator norm of a commutator $[\hat{X},\hat{Y}]$ for $\pm1$ observables satisfies $\|[\hat{X},\hat{Y}]\| \leq 2$, so $\|\hat{S}^2\| \leq 4 + 4 = 8$, giving $\|\hat{S}\| \leq 2\sqrt{2}$. The bound is tight: the singlet state with $45\circ$ settings achieves it exactly.

%% -----------------------------------------------------------
\subsection*{6. The intermediate bound and its philosophical meaning}

The existence of an intermediate bound---classical $\leq 2$, quantum $\leq 2\sqrt{2}$, logical maximum $= 4$---is not an accident, and its philosophical significance goes beyond confirming quantum mechanics.

\medskip

Popescu and Rohrlich noted in 1994 \cite{popescu_1994_pr} that one can define hypothetical theories, now called \textit{PR-boxes}, that respect the no-signalling principle (no faster-than-light communication) but saturate $|S| = 4$. Such theories are logically consistent and non-signalling, yet they would allow correlations far stronger than anything quantum mechanics predicts. The fact that nature chooses $2\sqrt{2}$, and not $4$, is therefore a \textit{non-trivial physical fact} that is not explained by no-signalling alone. Why this particular number?

\medskip

Several research programmes have tried to derive $2\sqrt{2}$ from information-theoretic principles---non-trivial communication complexity \cite{brassard_2006_limits}, information causality \cite{pawlowski_2009_information}, and local orthogonality \cite{fritz_2013_local}---but no single principle has been universally accepted as the explanation.

\medskip

For our purposes, this is directly relevant to the potentiality reading defended in Section~5. The Tsirelson bound shows that the joint potentialities encoded in the quantum state are not arbitrary: they obey a precise algebraic structure (the Hilbert space formalism, and specifically the constraint~(\ref{eq:S2}) on operator norms) that limits \textit{how strongly} the potentialities of two entangled systems can be correlated.

\medskip

The entanglement is not a free, unlimited non-local resource---it is a structured form of joint potentiality, constrained by the algebra of quantum observables. This is consistent with potentiality realism: the potentialities are \textit{real} (they exceed classical bounds), but they are \textit{structured} (they obey the Tsirelson bound rather than the logical maximum).

\medskip

The gap between $2\sqrt{2}$ and $4$ is, in this sense, a quantitative signature of the ontological structure of quantum potentiality---a structure that neither classical ontology nor pure logic would have predicted, and that remains only partially understood.
